%%%
%%% Radical "⼌"
%%%
\section*{Radical 13: ``⼌''}\addcontentsline{toc}{section}{Radical 13: ⼌}\addcontentsline{loh}{figure}{\#\#\#\# 13: ⼌}

%%%%%%%%%% 内 %%%%%%%%%%
\subsection*{内}\addcontentsline{loh}{figure}{内}

\begin{Entry}{内}{4}{⼌}
  \begin{Phonetics}{内}{nei4}[][HSK 3]
    \definition*{s.}{Sobrenome: Nei}
    \definition{s.}{dentro; interior; parte interna ou lateral |  (antes de um substantivo ou verbo na formação de uma palavra composta) interno | (depois de um substantivo para indicar lugar, tempo, escopo ou limites) dentro; em | coração; mente | esposa ou parentes da esposa}
  \antonymref{外}{wai4}
  \end{Phonetics}
\end{Entry}

\begin{Entry}{内心}{4,4}{⼌,⼼}
  \begin{Phonetics}{内心}{nei4xin1}[][HSK 3]
    \definition{s.}{coração; interior; íntimo do ser}
  \end{Phonetics}
\end{Entry}

\begin{Entry}{内外}{4,5}{⼌,⼣}
  \begin{Phonetics}{内外}{nei4wai4}[][HSK 6]
    \definition{s.}{dentro e fora; nacional e estrangeiro; interno e externo | ao redor; aproximadamente; número aproximado de exibição}
  \end{Phonetics}
\end{Entry}

\begin{Entry}{内向}{4,6}{⼌,⼝}
  \begin{Phonetics}{内向}{nei4xiang4}[][HSK 7-9]
    \definition{adj.}{(economia, etc.) voltado para o mercado interno; introvertido; descreve a personalidade de uma pessoa como sendo quieta e relutante em expressar seus pensamentos}
  \end{Phonetics}
\end{Entry}

\begin{Entry}{内在}{4,6}{⼌,⼟}
  \begin{Phonetics}{内在}{nei4zai4}[][HSK 5]
    \definition{adj.}{intrínseco; algo que existe em si mesmo, mas que não pode ser descoberto através da observação direta | interno; imanente; difícil de perceber}
  \end{Phonetics}
\end{Entry}

\begin{Entry}{内地}{4,6}{⼌,⼟}
  \begin{Phonetics}{内地}{nei4di4}[][HSK 6]
    \definition{s.}{interior; sertão | China continental (RPC excluindo Hong Kong e Macau, mas incluindo ilhas como Hainan) | Japão (usado em Taiwan durante a colonização japonesa)}
  \end{Phonetics}
\end{Entry}

\begin{Entry}{内存}{4,6}{⼌,⼦}
  \begin{Phonetics}{内存}{nei4cun2}[][HSK 7-9]
    \definition{s.}{RAM (\emph{random access memory}) | capacidade da RAM | memória do computador; memória interna, abreviação de 内存储器 | armazenamento interno}
  \seealsoref{内存储器}{nei4cun2chu3qi4}
  \seealsoref{随机存取存储器}{sui2ji1cun2qu3 cun2chu3qi4}
  \seealsoref{随机存取记忆体}{sui2ji1cun2qu3 ji4yi4ti3}
  \end{Phonetics}
\end{Entry}

\begin{Entry}{内存储器}{4,6,12,16}{⼌,⼦,⼈,⼝}
  \begin{Phonetics}{内存储器}{nei4cun2chu3qi4}
    \definition{s.}{memória interna}
  \end{Phonetics}
\end{Entry}

\begin{Entry}{内行}{4,6}{⼌,⾏}
  \begin{Phonetics}{内行}{nei4hang2}[][HSK 7-9]
    \definition{adj.}{habilidoso; proficiente; especialista em; amplo conhecimento e experiência em determinado assunto ou função}
    \definition{s.}{especialista; perito; craque; profissional}
  \end{Phonetics}
\end{Entry}

\begin{Entry}{内衣}{4,6}{⼌,⾐}
  \begin{Phonetics}{内衣}{nei4yi1}[][HSK 6]
    \definition[件,个]{s.}{roupa íntima}
  \end{Phonetics}
\end{Entry}

\begin{Entry}{内省}{4,9}{⼌,⽬}
  \begin{Phonetics}{内省}{nei4xing3}
    \definition{s.}{introspecção}
    \definition{v.}{refletir sobre si mesmo}
  \synonymref{反省}{fan3xing3}
  \synonymref{宫廷}{gong1ting2}
  \synonymref{皇宫}{huang2gong1}
  \antonymref{民间}{min2jian1}
  \end{Phonetics}
\end{Entry}

\begin{Entry}{内科}{4,9}{⼌,⽲}
  \begin{Phonetics}{内科}{nei4ke1}[][HSK 4]
    \definition{s.}{medicina geral; clínica geral; clínica médica}
  \end{Phonetics}
\end{Entry}

\begin{Entry}{内阁}{4,9}{⼌,⾨}
  \begin{Phonetics}{内阁}{nei4ge2}[][HSK 7-9]
    \definition{s.}{gabinete; em alguns países, o órgão executivo máximo é composto pelo primeiro"-ministro e vários membros do gabinete (ministros, chefes de gabinete ou outros funcionários de alto escalão)}
  \end{Phonetics}
\end{Entry}

\begin{Entry}{内容}{4,10}{⼌,⼧}
  \begin{Phonetics}{内容}{nei4rong2}[][HSK 3]
    \definition[份,个,项]{s.}{conteúdo; substância; a substância ou significado contido em algo}
  \end{Phonetics}
\end{Entry}

\begin{Entry}{内资}{4,10}{⼌,⾙}
  \begin{Phonetics}{内资}{nei4zi1}
    \definition{s.}{capital nacional; financiamento interno; investimento de fontes nacionais}
  \antonymref{外资}{wai4zi1}
  \end{Phonetics}
\end{Entry}

\begin{Entry}{内部}{4,10}{⼌,⾢}
  \begin{Phonetics}{内部}{nei4bu4}[][HSK 4]
    \definition{s.}{interior; dentro; interno; dentro de um determinado intervalo}
  \end{Phonetics}
\end{Entry}

\begin{Entry}{内涵}{4,11}{⼌,⽔}
  \begin{Phonetics}{内涵}{nei4han2}[][HSK 7-9]
    \definition{s.}{intenção; conotação; os atributos essenciais das coisas objetivas refletidos por um conceito}
  \end{Phonetics}
\end{Entry}

\begin{Entry}{内幕}{4,13}{⼌,⼱}
  \begin{Phonetics}{内幕}{nei4mu4}[][HSK 7-9]
    \definition{s.}{o que acontece nos bastidores; histórias internas}
  \end{Phonetics}
\end{Entry}

\begin{Entry}{内需}{4,14}{⼌,⾬}
  \begin{Phonetics}{内需}{nei4xu1}[][HSK 7-9]
    \definition{s.}{Economia: demanda interna}
  \antonymref{外需}{wai4xu1}
  \end{Phonetics}
\end{Entry}

\begin{Entry}{内燃机}{4,16,6}{⼌,⽕,⽊}
  \begin{Phonetics}{内燃机}{nei4ran2ji1}
    \definition{s.}{motor de combustão interna; um tipo de motor térmico, no qual o combustível queima no cilindro, produzindo gases em expansão que empurram o pistão, que por sua vez aciona a biela para girar o eixo do motor; os motores de combustão interna utilizam gasolina, etanol; diesel ou gás natural como combustível}
  \end{Phonetics}
\end{Entry}

%%%%%%%%%% 册 %%%%%%%%%%
\subsection*{册}\addcontentsline{loh}{figure}{册}

\begin{Entry}{册}{5}{⼌}
  \begin{Phonetics}{册}{ce4}[][HSK 5]
    \definition{clas.}{usado para cópias de livros}
    \definition{s.}{volume; livro | cópia; volume | ordem imperial para conferir um título}
    \definition{v.}{conferir um título}
  \end{Phonetics}
\end{Entry}

%%%%%%%%%% 再 %%%%%%%%%%
\subsection*{再}\addcontentsline{loh}{figure}{再}

\begin{Entry}{再}{6}{⼌}
  \begin{Phonetics}{再}{zai4}[][HSK 1]
    \definition{adv.}{mais uma vez; além disso; ainda mais; indica a repetição ou continuação de uma mesma ação ou comportamento; refere"-se principalmente a ações ou comportamentos não realizados ou contínuos | usado antes do adjetivo, indica intensificação, equivalente a 更 ou 更加 | (para uma ação adiada, precedida por uma expressão de tempo ou condição) então; somente então; depois de algo; indica que a ação ocorrerá após a conclusão de outra ação | além disso; indica um complemento, equivalente a 另外 ou 又 | próxima vez; indica que a ação ocorrerá após um determinado período de tempo | novamente; de novo}
  \seealsoref{重新}{chong2xin1}
  \seealsoref{更}{geng4}
  \seealsoref{更加}{geng4jia1}
  \seealsoref{还}{hai2}
  \seealsoref{另外}{ling4wai4}
  \seealsoref{又}{you4}
  \synonymref{重新}{chong2xin1}
  \synonymref{还}{hai2}
  \synonymref{又}{you4}
  \end{Phonetics}
\end{Entry}

\begin{Entry}{再三}{6,3}{⼌,⼀}
  \begin{Phonetics}{再三}{zai4san1}[][HSK 4]
    \definition{adv.}{repetidamente; repetidas vezes; de novo e de novo}
  \end{Phonetics}
\end{Entry}

\begin{Entry}{再也}{6,3}{⼌,⼄}
  \begin{Phonetics}{再也}{zai4ye3}[][HSK 5]
    \definition{adv.}{não mais; nunca mais; uma determinada situação ou ação nunca mais ocorrerá}
  \end{Phonetics}
\end{Entry}

\begin{Entry}{再不}{6,4}{⼌,⼀}
  \begin{Phonetics}{再不}{zai4 bu4}
    \definition{adv.}{ou então; ou}
  \end{Phonetics}
\end{Entry}

\begin{Entry}{再见}{6,4}{⼌,⾒}
  \begin{Phonetics}{再见}{zai4jian4}[][HSK 1]
    \definition{v.}{adeus; tchau; até logo; até mais; até mais tarde}
  \end{Phonetics}
\end{Entry}

\begin{Entry}{再发}{6,5}{⼌,⼜}
  \begin{Phonetics}{再发}{zai4fa1}
    \definition{v.}{reenviar; reeditar}
  \end{Phonetics}
\end{Entry}

\begin{Entry}{再生}{6,5}{⼌,⽣}
  \begin{Phonetics}{再生}{zai4sheng1}[][HSK 6]
    \definition{v.}{reviver; ressuscitar; ressuscitar dos mortos | reproduzir; regenerar | reprocessar; reciclar; regenerar; processar um determinado produto residual para restaurar seu desempenho original e transformá"-lo em um novo produto}
  \end{Phonetics}
\end{Entry}

\begin{Entry}{再次}{6,6}{⼌,⽋}
  \begin{Phonetics}{再次}{zai4ci4}[][HSK 5]
    \definition{adv.}{mais uma vez; uma segunda vez; outra vez}
  \end{Phonetics}
\end{Entry}

\begin{Entry}{再审}{6,8}{⼌,⼧}
  \begin{Phonetics}{再审}{zai4shen3}
    \definition{s.}{revisão (de um caso, orçamento, artigo, etc.) | novo julgamento; reabertura de um caso; nova audiência; revisão judicial}
    \definition{v.}{ter um novo julgamento | examinar novamente; reexaminar}
  \end{Phonetics}
\end{Entry}

\begin{Entry}{再现}{6,8}{⼌,⾒}
  \begin{Phonetics}{再现}{zai4xian4}[][HSK 7-9]
    \definition{v.}{(um evento passado) reaparecer; ser reproduzido}
  \synonymref{表现}{biao3xian4}
  \synonymref{重现}{chong2xian4}
  \synonymref{体现}{ti3xian4}
  \end{Phonetics}
\end{Entry}

\begin{Entry}{再者}{6,8}{⼌,⽼}
  \begin{Phonetics}{再者}{zai4zhe3}
    \definition{adv./conj.}{além disso; ademais; além de}
  \synonymref{此外}{ci3wai4}
  \synonymref{另外}{ling4wai4}
  \end{Phonetics}
\end{Entry}

\begin{Entry}{再育}{6,8}{⼌,⾁}
  \begin{Phonetics}{再育}{zai4yu4}
    \definition{v.}{aumentar | multiplicar | proliferar}
  \end{Phonetics}
\end{Entry}

\begin{Entry}{再临}{6,9}{⼌,⼁}
  \begin{Phonetics}{再临}{zai4lin2}
    \definition{v.}{vir de novo; voltar}
  \end{Phonetics}
\end{Entry}

\begin{Entry}{再度}{6,9}{⼌,⼴}
  \begin{Phonetics}{再度}{zai4du4}[][HSK 7-9]
    \definition{adv.}{mais uma vez; uma segunda vez; novamente}
  \synonymref{再次}{zai4ci4}
  \end{Phonetics}
\end{Entry}

\begin{Entry}{再说}{6,9}{⼌,⾔}
  \begin{Phonetics}{再说}{zai4shuo1}[][HSK 6]
    \definition{conj.}{além disso; o que é mais; indicando uma razão adicional, equivalente a 况且}
    \definition{v.}{adiar para mais tarde; deixar para processamento ou consideração posterior}
  \seealsoref{况且}{kuang4qie3}
  \end{Phonetics}
\end{Entry}

\begin{Entry}{再读}{6,10}{⼌,⾔}
  \begin{Phonetics}{再读}{zai4 du2}
    \definition{v.}{ler novamente | revisar (uma lição etc.)}
  \end{Phonetics}
\end{Entry}

%%%%%%%%%% 罔 %%%%%%%%%%
\subsection*{罔}\addcontentsline{loh}{figure}{罔}

\begin{Entry}{罔}{8}{⼌}
  \begin{Phonetics}{罔}{wang3}
    \definition{adj.}{Literário: não disponível}
    \definition{adv.}{Literário: não; nenhum}
    \definition{s.}{confuso e ignorante}
    \definition{v.}{Literário: enganar; trapacear; encobrir}
  \synonymref{惑}{huo4}
  \synonymref{没}{mei2}
  \synonymref{蒙}{meng1}
  \synonymref{迷}{mi2}
  \antonymref{有}{you3}
  \end{Phonetics}
\end{Entry}

%%%%% EOF %%%%%

